\documentclass[10pt,a4paper]{amsart}
\usepackage{amsmath,amssymb,amsthm,graphicx}
% Retained minimal preamble for the Stacks Project chapter more-algebra.tex (More on Algebra).
% The chapter file pinned-original-source.tex carries only \input{preamble} on its first line, so no
% preamble text lives inside the chapter itself.  The Stacks Project uses one single shared preamble for
% every chapter; the snapshot stages it identically as sources/stacks-chow/preamble.tex and
% sources/stacks-etale-cohomology/preamble.tex.  The block below is copied mechanically, line for line,
% from sources/stacks-chow/preamble.tex lines 14-18 (the Xy-pic setup) and 154-217 (the native theorem
% declarations together with the mathematical macros).  Nothing was added, removed or redefined; the
% wrapper class and the amsmath/amssymb/amsthm/graphicx base packages are supplied by the assembler.
% Chapter: more-algebra.tex -- More on Algebra.
% For commutative diagrams we use Xy-pic
\usepackage[all]{xy}

% We use 2cell for 2-commutative diagrams.
\xyoption{2cell}

% Theorem environments.
%
\theoremstyle{plain}
\newtheorem{theorem}[subsection]{Theorem}
\newtheorem{proposition}[subsection]{Proposition}
\newtheorem{lemma}[subsection]{Lemma}

\theoremstyle{definition}
\newtheorem{definition}[subsection]{Definition}
\newtheorem{example}[subsection]{Example}
\newtheorem{exercise}[subsection]{Exercise}
\newtheorem{situation}[subsection]{Situation}

\theoremstyle{remark}
\newtheorem{remark}[subsection]{Remark}
\newtheorem{remarks}[subsection]{Remarks}

\numberwithin{equation}{subsection}

% Macros
%
\def\lim{\mathop{\mathrm{lim}}\nolimits}
\def\colim{\mathop{\mathrm{colim}}\nolimits}
\def\Spec{\mathop{\mathrm{Spec}}}
\def\Hom{\mathop{\mathrm{Hom}}\nolimits}
\def\Ext{\mathop{\mathrm{Ext}}\nolimits}
\def\SheafHom{\mathop{\mathcal{H}\!\mathit{om}}\nolimits}
\def\SheafExt{\mathop{\mathcal{E}\!\mathit{xt}}\nolimits}
\def\Sch{\mathit{Sch}}
\def\Mor{\mathop{\mathrm{Mor}}\nolimits}
\def\Ob{\mathop{\mathrm{Ob}}\nolimits}
\def\Sh{\mathop{\mathit{Sh}}\nolimits}
\def\NL{\mathop{N\!L}\nolimits}
\def\CH{\mathop{\mathrm{CH}}\nolimits}
\def\proetale{{pro\text{-}\acute{e}tale}}
\def\etale{{\acute{e}tale}}
\def\QCoh{\mathit{QCoh}}
\def\Ker{\mathop{\mathrm{Ker}}}
\def\Im{\mathop{\mathrm{Im}}}
\def\Coker{\mathop{\mathrm{Coker}}}
\def\Coim{\mathop{\mathrm{Coim}}}

% Boxtimes
%
\DeclareMathSymbol{\boxtimes}{\mathbin}{AMSa}{"02}

%
% Macros for moduli stacks/spaces
%
\def\QCohstack{\mathcal{QC}\!\mathit{oh}}
\def\Cohstack{\mathcal{C}\!\mathit{oh}}
\def\Spacesstack{\mathcal{S}\!\mathit{paces}}
\def\Quotfunctor{\mathrm{Quot}}
\def\Hilbfunctor{\mathrm{Hilb}}
\def\Curvesstack{\mathcal{C}\!\mathit{urves}}
\def\Polarizedstack{\mathcal{P}\!\mathit{olarized}}
\def\Complexesstack{\mathcal{C}\!\mathit{omplexes}}
% \Pic is the operator that assigns to X its picard group, usage \Pic(X)
% \Picardstack_{X/B} denotes the Picard stack of X over B
% \Picardfunctor_{X/B} denotes the Picard functor of X over B
\def\Pic{\mathop{\mathrm{Pic}}\nolimits}
\def\Picardstack{\mathcal{P}\!\mathit{ic}}
\def\Picardfunctor{\mathrm{Pic}}
\def\Deformationcategory{\mathcal{D}\!\mathit{ef}}
\usepackage{cleveref}
\crefname{originallemma}{original lemma}{original lemmas}
\Crefname{originallemma}{Original lemma}{Original lemmas}
\makeatletter
\newlabel{lemma-H0-inverse}{{05EM}{source}{}{sourcepointer-1}{}}
\newlabel{lemma-H0-inverse@cref}{{[originallemma][0][]05EM}{[1][0][]source}{}{}{}}
\makeatother
\crefname{originallemma}{original lemma}{original lemmas}
\Crefname{originallemma}{Original lemma}{Original lemmas}
\makeatletter
\newlabel{lemma-base-change-glue}{{05EQ}{source}{}{sourcepointer-2}{}}
\newlabel{lemma-base-change-glue@cref}{{[originallemma][0][]05EQ}{[1][0][]source}{}{}{}}
\makeatother
\crefname{originallemma}{original lemma}{original lemmas}
\Crefname{originallemma}{Original lemma}{Original lemmas}
\makeatletter
\newlabel{lemma-equivalence-I-unit}{{05EP}{source}{}{sourcepointer-3}{}}
\newlabel{lemma-equivalence-I-unit@cref}{{[originallemma][0][]05EP}{[1][0][]source}{}{}{}}
\makeatother
\crefname{originallemma}{original lemma}{original lemmas}
\Crefname{originallemma}{Original lemma}{Original lemmas}
\makeatletter
\newlabel{lemma-neighbourhood-extensions}{{05EG}{source}{}{sourcepointer-4}{}}
\newlabel{lemma-neighbourhood-extensions@cref}{{[originallemma][0][]05EG}{[1][0][]source}{}{}{}}
\makeatother
\title{The functors Can and H\textasciicircum{}0 are quasi-inverse equivalences for the glued category of a flat cover with a finitely generated ideal}
\author{The Stacks Project Authors}
\date{stacks-project:more-algebra.tex:local-snapshot-20261012; GNU FDL 1.2 or later}
\begin{document}
\maketitle
\setcounter{subsection}{1}
\noindent\textit{Editorial note.} All delivered mathematical text is the author's own, copied from the pinned Stacks Project chapter file more-algebra.tex; no statement or proof step was rewritten, reordered or removed. The counted statement is the labeled proposition environment at source lines 24812-24821 (canonical tag 05ER) and the counted proof is the proof environment at source lines 24823-24879, which immediately follows it and is its own proof.\par\medskip
\noindent\textit{Editorial note.} The proof displays the comparison morphism and, later, a commutative diagram with exact rows in the author's own Xy-pic markup, both of which are compiled from that markup with the shared preamble.\par\medskip
\noindent\textit{Editorial note.} The proof ends with the parenthetical remark that a posteriori the module Q vanishes; the remark is the author's own and is preserved verbatim inside the counted span.\par\medskip
\noindent\textit{Editorial note.} The chapter file pinned-original-source.tex carries only an input of the shared preamble on its first line, so no preamble text lives inside the chapter. The retained file sources/stacks-more-algebra/retained-preamble.tex is a mechanical, line-for-line copy of the project's single shared preamble as staged in the snapshot (Xy-pic setup plus the native theorem declarations and mathematical macros). Nothing was added, removed or redefined, and the excerpt serves only as preamble, never as a span coordinate.\par\medskip
\noindent\textit{Editorial note.} The source URL is the source-level Stacks Project URL, as required when more than one card is produced from one source; the per-claim tag is carried by claim\_key and by this note. The counted statement and proof cite no bibliography entry, so this card needs no bibliography.\par\medskip
\noindent\textit{Allowed standard prerequisites (editorial).}\begin{itemize}
\item The composite of the functor Can with H\textasciicircum{}0 is isomorphic to the identity functor on the category of modules. used at proof lines 24823-24826. Disposes of one of the two composites, so that only the composite starting from a glued object has to be analysed.
\item After inverting one of the generators the comparison morphism becomes an isomorphism, and its kernel and cokernel are then systems concentrated in the first component. used at proof lines 24841-24846. Reduces the analysis of the kernel and cokernel of the comparison morphism to systems of the special form with only the S-component present.
\item A system of the form with only the first component present has the property that H\textasciicircum{}0 of it is that component, and H\textasciicircum{}0 is left exact, so the exactness computation gives that the kernel vanishes. used at proof lines 24846-24849. Shows that the comparison morphism is injective, which is the first half of the proof that it is an isomorphism.
\item For an I-power torsion module the canonical map from the tensor product with S is an isomorphism, and extensions of the module by the image of H\textasciicircum{}0 can be realized inside the glued category. used at proof lines 24856-24860. Produces the commutative diagram with exact rows whose middle term gives a module E with Can(E) isomorphic to the glued object.
\end{itemize}\medskip
\ifdefined\hypertarget\hypertarget{sourcepointer-1}{}\fi
\noindent\textit{Source reference.} Original lemma 05EM: sources/stacks-more-algebra/pinned-original-source.tex, source lines 24729--24729. The reference records that one of the two composites is the identity functor, used at proof lines 24824-24826; the analysis of the other composite, which is the content of the proof, is inside the counted span.\par\medskip
\ifdefined\hypertarget\hypertarget{sourcepointer-2}{}\fi
\noindent\textit{Source reference.} Original lemma 05EQ: sources/stacks-more-algebra/pinned-original-source.tex, source lines 24792--24792. The reference is combined with the next one at proof lines 24842-24844 to show that the comparison morphism becomes an isomorphism after inverting a generator; the conclusions drawn from that localization are inside the counted span.\par\medskip
\ifdefined\hypertarget\hypertarget{sourcepointer-3}{}\fi
\noindent\textit{Source reference.} Original lemma 05EP: sources/stacks-more-algebra/pinned-original-source.tex, source lines 24768--24768. The reference is the second of the two lemmas applied at proof lines 24842-24844 to localize the comparison morphism; the ensuing injectivity and exactness arguments are inside the counted span.\par\medskip
\ifdefined\hypertarget\hypertarget{sourcepointer-4}{}\fi
\noindent\textit{Source reference.} Original lemma 05EG: sources/stacks-more-algebra/pinned-original-source.tex, source lines 24477--24477. The reference supplies the existence of the commutative diagram with exact rows used at proof lines 24858-24874 to produce the module E; the argument that this E maps under Can to the given glued object is inside the counted span.\par\medskip
\begin{proposition}
\label{proposition-equivalence}
Assume $\varphi : R \to S$ is a flat ring map and
$I = (f_1, \ldots, f_t) \subset R$ is an ideal such that
$R/I \to S/IS$ is an isomorphism. Then $\text{Can}$ and
$H^0$ are quasi-inverse equivalences of categories
$$
\text{Mod}_R = \text{Glue}(R \to S, f_1, \ldots, f_t)
$$
\end{proposition}

\begin{proof}
We have already seen that $H^0 \circ \text{Can}$ is isomorphic to the
identity functor, see
Lemma \ref{lemma-H0-inverse}.
Consider an object $\mathbf{M} = (M', M_i, \alpha_i, \alpha_{ij})$
of $\text{Glue}(R \to S, f_1, \ldots, f_t)$.
We get a natural morphism
$$
\Psi :
(H^0(\mathbf{M}) \otimes_R S, H^0(\mathbf{M})_{f_i},
\text{can}_i, \text{can}_{ij})
\longrightarrow
(M', M_i, \alpha_i, \alpha_{ij}).
$$
Namely, by definition $H^0(\mathbf{M})$ comes equipped with compatible
$R$-module maps $H^0(\mathbf{M}) \to M'$ and $H^0(\mathbf{M}) \to M_i$.
We have to show that this map is an isomorphism.

\medskip\noindent
Pick an index $i$ and set $R' = R_{f_i}$. Combining
Lemmas \ref{lemma-base-change-glue} and \ref{lemma-equivalence-I-unit}
we see that $\Psi \otimes_R R'$ is an isomorphism.
Hence the kernel, resp.\ cokernel of $\Psi$ is a system of the form
$(K, 0, 0, 0)$, resp.\ $(Q, 0, 0, 0)$. Note that
$H^0((K, 0, 0, 0)) = K$, that $H^0$ is left exact, and that by
construction $H^0(\Psi)$ is bijective. Hence we see $K = 0$, i.e.,
the kernel of $\Psi$ is zero.

\medskip\noindent
The conclusion of the above is that we obtain a short exact sequence
$$
0 \to H^0(\mathbf{M}) \otimes_R S \to M' \to Q \to 0
$$
and that $M_i = H^0(\mathbf{M})_{f_i}$.
Note that we may think of $Q$ as an $R$-module which is $I$-power
torsion so that $Q = Q \otimes_R S$. By
Lemma \ref{lemma-neighbourhood-extensions}
we see that there exists a commutative diagram
$$
\xymatrix{
0 \ar[r] &
H^0(\mathbf{M}) \ar[r] \ar[d] &
E \ar[r] \ar[d] &
Q \ar[r] \ar[d] &
0 \\
0 \ar[r] &
H^0(\mathbf{M}) \otimes_R S \ar[r] &
M' \ar[r] &
Q \ar[r] &
0
}
$$
with exact rows. This clearly determines an isomorphism
$\text{Can}(E) \to (M', M_i, \alpha_i, \alpha_{ij})$
in the category $\text{Glue}(R \to S, f_1, \ldots, f_t)$
and we win. (Of course, a posteriori we have $Q = 0$.)
\end{proof}

\end{document}
